% !TEX root = ../main.tex
\lecture{9}{2026-11-16}{Akteure des schweizerischen Finanzplatzes}

\subsection{Finanzmarktaufsicht (FINMA)}

\subsubsection{Aufgaben}

\begin{itemize}
    \item Erteilt Bewilligung für Bankentätigkeit,
    Versicherungen und Effektenhändler
    \item Überwachung und Durchsetzung des Aufsichtsrechts
    \item Regulationstätigkeit (erlässt Verordnungen und
    Rundschreiben)
    \item Zusammengefasst: Sicherstellen dass Akteure am
    Finanzplatz legitim und Gesetze folgen, sodass Stabilität
    bewahrt wird
\end{itemize}

\subsection{Börse SIX Group}
\begin{definition}
    [Swiss Infrastructure and Exchange] ist eine
    privatrechtliche AG von verschiedenen gross Banken besitzt. Entwickelt die
    Finanzmarktinfrastruktur und führt Finanzdienstleistungen
\end{definition}

\subsubsection{Services}
\begin{enumerate}
    \item Exchange
    \item Securities Services
    \item Financial Information: Referenz und Kursdaten
    \item Banking Services: bargeldlosen instant Zahlungen,
    Interbank Zahlungsverkehr 
\end{enumerate}

\subsection{Weitere Akteure}
\begin{itemize}
    \item Einlagensicherung (esisuisse)
    \begin{itemize}
        \item Sichert privilegierte Kudeneinlagen bei Banken
        \item Um Bank-Runs bis 100k zu minimieren
    \end{itemize}
    \item Bankenombudsman: Neutrale und kostenlose Anlaufstelle für Personen
    ― Bei Beschwerden von Kund:innen gegen Banken in der Schweiz
    ― Bei nachrichtenlosen Vermögenswerten (schlafende Konti)
    ― Bank für Internationalen Zahlungsausgleich (BIZ)
    ― Kooperation zur Förderung Finanzstabilität, Wissensaustausch, Erhöhung Aufsichtsqualität
    → erarbeitet Standards zu Bankenaufsicht «Basel I, II, III
\end{itemize}


\section{Commercial/Retail Banking}


\subsection{Pfand}
Besicherung der Kredite gegen Stellung von Sicherheiten

\begin{itemize}
    \item Zweck
    \begin{itemize}
        \item Verstärkung der Anreize der Rückzahlung
        \item Verminderung eines allfälligen Verlusts
        \item Signaling
    \end{itemize}
    \item Gängige Sicherheiten
    \begin{itemize}
        \item Immobilien
        \item Wertschriften
        \item Forderungen und Warenlager 
    \end{itemize}
\end{itemize}

\subsection{Hypothekargeschäft}

\subsubsection{Kriterien für Hypothekarvergabe}

\begin{itemize}
    \item Tragbarkeit
    \begin{itemize}
        \item Messung ob Haushalt die langfristigen Kosten der
        Hypothek aus dem laufenden Einkommen tragen kann
        \item Tragbarkeitsberechnung
        \begin{itemize}
            \item jährliche Kosten für Hypothekarzinsen (mit
            kalkulatorischem Zins 5\%), Amortisation und
            Nebenkosten(1\% des Immowerts) nicht > 35\% des Bruttoeinkommens
        \end{itemize}
    \end{itemize}
    \item Belehnung
    \begin{itemize}
        \item Verhältnis zum mindestkaufpreis der liegenschaft
        und Hypothek, wobei Käufer mind. 20\% EK einbringen muss
    \end{itemize}
\end{itemize}



